\clearpage
\begin{summarybox}
\textbf{\textcolor{CSummary}{这条结论是为了解决什么？}}

题目不告诉 $f$ 的具体式子，只告诉一部分定义域信息。
这时先区分\textbf{$x$ 是什么数}、\textbf{括号里真正送给 $f$ 的又是什么数}，
便能直接算出允许范围。

\medskip
\textbf{\textcolor{CSummary}{考场两步法}}
\begin{description}[style=nextline,leftmargin=2.8em,labelsep=0.5em,itemsep=0.4em]
 \item[\textbf{已知 $f(t)$，求 $f[g(x)]$：}]
 让 $g(x)$ 有意义，并且落入 $f$ 原来允许的输入范围；\textbf{解出 $x$}。
 \item[\textbf{已知 $f(ax+b)$，求 $f(t)$（$a\ne0$）：}]
 让 $x$ 在题目给出的\textbf{完整定义域}内变化；\textbf{算出 $ax+b$ 的范围}。
\end{description}

\textbf{\textcolor{CSummary}{三个必须检查的细节}}
\begin{enumerate}[leftmargin=*,itemsep=0.25em]
 \item 定义域始终回答“当前函数的\textbf{自变量 $x$} 能取什么值”。
 \item 解不等式要保留端点是否包含；括号里的根号、分母本身也须有意义。
 \item \textbf{反推完整定义域前，先检查括号里的式子能否取遍所有实数。}
 $x^2$、$|x|$ 等不能取负数，不能照搬 $ax+b$ 的逆推结论。
\end{enumerate}

\medskip
\textbf{\textcolor{CSummary}{一句话记住}}

\textbf{正着求，看送进去的数合不合格；反着求，看送进去的数能不能把所有情况都覆盖。}
\end{summarybox}
